\subsection*{13.6 Application: Martingale and Stopping Time}

The term "martingale" originated from French, where it referred to a half belt used

to secure a horse or dog's neck. In probability theory, a martingale is a stochastic

process that models a fair game in which the expected gain of a gambler is zero,

assuming that the gambler can only make decisions in a causal manner. This

probability model has found wide applications in finance.

We encode the information available from the history in the form of a filtration.

(cf. Example 2.2.1)

\begin{defbox}{13.6}
\end{defbox}

A filtration is a sequence of increasing \sigma-subfields in \mathcal{F}, denoted as .(\mathcal{F}n)\infty

n=0,

such that

\mathcal{F}0 \subseteq \mathcal{F}1 \subseteq \mathcal{F}2 \subseteq\cdot\cdot\cdot .

A sequence of random variables .(Xn)\infty

n=0 is said to be adapted to filtration

(\mathcal{F}n)\infty

n=0 if. Xn is\mathcal{F}n-measurable for all n.

Another way to describe a filtration is to regard the index n as representing

discrete time. At each time n, the filtration captures the information available up

to that point in time. Typically, the first \sigma-field \mathcal{F}0 in the filtration is the trivial

\sigma-field.{\emptyset,\Omega }, ie., at time n = 0, only the trivial events \emptyset and. \Omega are measurable.

The defining property of a filtration is that if something is measurable at time n,

then it continues to be measurable at any future time n' >n As time advances, we

have more available information, and the \sigma-algebras in the filtration become larger.

\textbf{Example 13.6.1 (Natural Filtration)} \\

Suppose Y1,Y 2,.is a sequence of random variables defined on probability space .(\Omega, \mathcal{F},P) .

Let \mathcal{F}0 \coloneqq {\emptyset,\Omega } and \mathcal{F}n \coloneqq \sigma(Y1,Y 2,...,Y n),f o r n = 1, 2, 3,... The n .(\mathcal{F}n)n\geq0 is a

filtration. If fn(y1,...,y n) is any \mathcal{B}(Rn)-measurable function defined on Rn, then the sequence

(fn(Y1,...,Y n))\infty

i=1 is adapted to .(\mathcal{F}n)n\geq0. This means that fn(Y1,Y 2,...,Y n) depends only on

the information available up to and including time n. This filtration is called the natural filtration

for the sequence .(Yi)i\geq1.

\begin{defbox}{13.7}
\end{defbox}

A sequence of random variables .(Xn)\infty

n=0 is a martingale relative to filtration

(\mathcal{F}n)\infty

n=0 if it satisfies the following three conditions:

1. E[\vertXn\vert] < \infty for all n, ie.,Xn \in L1(P) .

2. .(Xn)\infty

n=0 is adapted to .(\mathcal{F}n)\infty

n=0.

3. E[Xn+1\vert\mathcal{F}n]= Xn as. forn \geq 0.

The third condition captures the idea that the expected value of the next

observation, given the current and all past observations, is equal to the current

observation.

When the filtration is defined by a sequence of random variables .(Yn)n\geq0, then

random variable Xn is a measurable function of .(Y0,Y 1,...,Y n) (Theorem 12.6).

We can express the last condition in Definition 13.7 as E[Xn+1\vertY0,Y 1,...,Y n]=

Xn.

Martingale can be used to model a variety of stochastic processes, including a

betting system, where Xn represents the amount of money in the gambler's pocket

at time n The essence of a martingale is that, regardless of the betting strategy

adopted, the expected value of Xn+1 in the next step, given the history up to time

n, is always equal to the current value Xn. In other words, a martingale represents a

fair game.

\begin{thmbox}{13.15}
\end{thmbox}

In Definition 13.7, we can replace the third condition by

E[Xm\vert\mathcal{F}n]= Xn as. for allm \geq n. (13.18)

\textit{Proof} To show that (13.18) implies the third condition in Definition 13.7, we can

take m= n + 1 To prove the reverse direction, suppose m\geq n + 1 B yt h e

tower property (Theorem 13.8), we have E[Xm\vert\mathcal{F}n]= E[E[Xm\vert\mathcal{F}m- 1]\vert\mathcal{F}n]=

E[Xm- 1\vert\mathcal{F}n] as. Repeating this argument, we obtain E[Xm\vert\mathcal{F}n] = \cdot\cdot\cdot =

E[Xn+1\vert\mathcal{F}n]= Xn as. \hfill $\square$

In a martingale, the unconditioned expectation E[Xn] remains constant for all n .

\begin{thmbox}{13.16}
\end{thmbox}

If.(Xn)\infty

n=0 is a martingale relative to some filtration .(\mathcal{F}n)\infty

n=0, then. E[Xn]=

E[X0] for alln\geq 1 .

\textit{Proof} Take expectation on both sides of (13.18)\hfill $\square$

\textbf{Example 13.6.2 (Random Walk as Martingale)} \\

Consider a sequence.(Yn)n\geq1 of iid. random variables with zero mean, and define the partial sums

Sn \coloneqqY 1+Y2+\cdot\cdot\cdot+ Yn forn\geq 1 ,w i t hS0 \coloneqq0 Let\mathcal{F}n \coloneqq\sigma(Y 1,...,Y n) be the natural filtration for

the sequence.(Yn)n\geq1 and set\mathcal{F}0 \coloneqq{\emptyset ,\Omega }It can be verified that S0,S 1,S 2,.form a martingale

relative to the filtration .(\mathcal{F}n)n\geq0.

To see this, we first note that Sn is a function of Y1,Y 2,...,Y n and is thus \mathcal{F}n-measurable.

Moreover, since. Yn has zero mean, we have

E[Sn+1\vert\mathcal{F}n]= E[Y1 + Y2 +\cdot\cdot\cdot+ Yn\vert\mathcal{F}n]+ E[Yn+1\vert\mathcal{F}n]

=E [Y1 + Y2 +\cdot\cdot\cdot+ Yn\vert\mathcal{F}n]+ E[Yn+1]

=Y 1 + Y2 +\cdot\cdot\cdot+ Yn + 0 = Sn,

where the last equality follows from the independence ofYn+1 from. \mathcal{F}n and the zero mean ofYn+1.

Finally, since\vertSn\vert\leq \sumn

i=1 \vertYi\vert,we haveE[\vertSn\vert] \leq \sumn

i=1 E[\vertYi\vert] < \infty for each n.

\textbf{Example 13.6.3 (Doob's Martingale)} \\

SupposeY1,Y 2,.is a sequence of random variables, and let X be an integrable random variable

defined on the same probability space. Define the natural filtration .(\mathcal{F}n)n\geq0 by \mathcal{F}0 ={ \emptyset,\Omega } and

\mathcal{F}n = \sigma(Y1,...,Y n),f o rn \geq 1LetZi \coloneqqE[X\vert\mathcal{F}i],f o ri = 0, 1, 2,... .

We note that Zn = E[X\vert\mathcal{F}n] is \mathcal{F}n-measurable by the definition of conditional expectation.

By the triangle inequality of conditioned expectation, we have

E[\vertZn\vert] = E[\vertE[X\vert\mathcal{F}n]\vert] \leq E[E[\vertX\vert\vert\mathcal{F}n]] = E[\vertX\vert] < \infty .

Finally, forn \geq 1, we can apply the tower property to obtain

E[Zn\vert\mathcal{F}n- 1]= E[E[X\vert\mathcal{F}n]\vert\mathcal{F}n- 1]= E[X\vert\mathcal{F}n- 1]= Zn- 1.

It proves that the sequence .(Zn)\infty

n=0 is a martingale sequence relative to .(\mathcal{F}n)n\geq0.

The stopping time T is a device that terminates a martingale. It is a random

variable such that the event .{T( )= n} can be determined by the information up to

time n. For example, a constant T is always a stopping time, but the first time that

a gambling loses ten times in a row is not, since it depends on future outcomes that

are not known at time n.

\begin{defbox}{13.8}
\end{defbox}

Given a filtration.(\mathcal{F}n)\infty

n=0 and a random variable

T : \Omega \to{ 0, 1, 2,..., }\cup{ \infty }

defined on .(\Omega, \mathcal{F},P) , we say that T is a stopping time relative to the filtration

(\mathcal{F}n)\infty

n=0 if the event.{T \leq n} is\mathcal{F}n-measurable for all n.

If.(Xn)\infty

n=0 is a sequence of random variables and T is a stopping time relative

to.(\sigma(X0,...,X n))n\geq0, we define a random variable XT by

XT () \coloneqqXT( )().

The random variable XT represents the value of the sequence .(Xn)n\geq0 at the

stopping time T , which is the time at which the martingale is terminated. Formally,

for each \in \Omega , we define XT () \coloneqqXT( )(), where T( )is the time at which

the stopping rule is applied to the realization . .

\textbf{Example 13.6.4 (Example of Stopping Time)} \\

Let.(Xn)\infty

n=0 be a sequence of random variables and let .(\mathcal{F}n)\infty

n=0 be the filtration defined by \mathcal{F}n =

\sigma(X0,X 1,...,X n)Define a random variable T as the first time n that the values of Xn exceed

100. Mathematically, we write

T( )=

{

min{n : Xn() \geq 100} if Xn() \geq 100 for some n

\infty otherwise.

When the sequence.(Xn())n\geq1 never exceeds the threshold 100, we define T( )to be infinity for

this realization. The random variable T is a stopping time relative to filtration .(\mathcal{F}n)\infty

n=0 because

the event.{ : T( )\leq n} can be written as the union of .{ : Xi() \geq 100},f o ri = 0, 1,...,n ,

which are all\mathcal{F}n-measurable.

\begin{defbox}{13.9}
\end{defbox}

Given a martingale .(Xn)n\geq0 and a stopping time T ,a stopped process . (XT

n )n\geq0

is obtained from .(Xn)n\geq0 by forcing the values with indices larger than T to be

XT That is, XT

n () \coloneqqXT n().

For allm \geq T , the random variable. XT

m stops updating and is equal to XT In the

definition, we used the notation a b \coloneqqmin(a, b).

\begin{thmbox}{13.17}
\end{thmbox}

The stopped process .(XT

n )n\geq0 is a martingale relative to the filtration

(\mathcal{F}n)n\geq0, andE[XT

n ]= E[X0] for alln \geq 0.

\textit{Proof} The expectationE[XT

n ] is less than or equal to maxk E[Xk], with maximum

taken over k = 0, 1,...,n Since E[Xk] is finite for all k, the maximum is also

finite.

For n \geq 0,t h e n-th random variable XT

n in the stopped process is a function of

X0,X 1,...,X n and T Since T is a stopping time, XT

n is\mathcal{F}n-measurable.

We write XT

n+1 as XT

n + (Xn+1 - Xn)1{T> n}The conditional expectation of

XT

n+1 given. \mathcal{F}n is

E[XT

n+1\vert\mathcal{F}n]= XT

n + 1{T> n}E[Xn+1 - Xn\vert\mathcal{F}n]= XT

n + 1{T> n} \cdot 0 = XT

n .

This proves that XT

n is a martingale relative to .(\mathcal{F}n)n\geq0.

Because.(XT

n )n\geq1 is a martingale, the last statement about the expectation. E[XT

n ]

follows from Theorem 13.16. \hfill $\square$

In contrast to the previous theorem, we are not just interested in the expected

value of the stopped process at a particular time, but also interested in the expected

value of the martingale sequence when it stops. The following theorem is called

the martingale stopping theorem, which is also known as Doob's optional stopping

theorem. Under the gambling scenario, the first condition in the theorem means

that a gambler has to leave the casino within a fixed duration. The second one

requires that the game must end when the gambler lose all of his fortune or win

a predetermined amount of money. The third one models the scenario in which the

increments are uniformly bounded.

\begin{thmbox}{13.18 (Martingale Stopping Theorem)}
\end{thmbox}

Let .(Xn)\infty

n=0 be a martingale and T be a stopping time relative to filtration

(\mathcal{F}n)\infty

n=0. Then

E[XT ]= E[X0]

if one of the following thee conditions hold:

(i) T is bounded by a constant C almost surely.

(ii) T is almost surely finite, and there exists a constant K such that \vertXn\vert\leq K

for all n .

(iii) E[T] < \infty , and there exists a constant c such that \vertXn -X n- 1\vert\leq c for

all n .

\textit{Proof} Under conditions (i), (ii), or (iii), the stopping time T is finite with

probability 1. Therefore, XT is well-defined almost surely and XT n \to X T as.

asn\to\infty .

(i) The martingale is stopped before time C as. We have XT =X C as. Hence,

E[XT ]= E[XC]= E[X0].

(ii) Since \vertXT n\vert is bounded by a constant for all n , by applying the dominated

convergence theorem, we obtain

limn\to\infty E[XT n]= E[Xlimn(T n)]= E[XT ].

In the second equality, we have used the assumption that T is finite as. and

thus limn(T n) = T as. From Theorem 13.17, we know that E[XT n]=

E[XT

n ]= E[X0] for all n Therefore,E[XT ]= E[X0].

(iii) SupposeE[T] is finite, and there exists a constant c such that. \vertXn -Xn- 1\vert\leq c

for all n We writeXT n as a telescoping sum

XT n =X 0 +

T n\sum

k=1

(Xk -X k- 1).

The number of terms is random but is no larger than n with probability 1.

Hence, we are dealing with finitely many additions in the calculation of XT n.

Using the bounded difference assumption, we obtain

\vertXT n\vert\leq\vert X0\vert+

T n\sum

k=1

\vertXk -X k- 1\vert\leq\vert X0\vert+ c\cdot (T n) \leq\vert X0\vert+ cT.

The right-hand side is integrable because both X0 and T are integrable. By

the dominated convergence theorem, we conclude that limn\to\infty E[XT n]=

E[XT ]Since E[XT n]= E[X0] for all n by Theorem 13.17, we obtain

E[XT ]= E[X0]\hfill $\square$

The next example is commonly known as the ABRACADABRA problem [ 10].

\textbf{Example 13.6.5 (Martingale Patterns and Gambling Team)} \\

Suppose a monkey type randomly on a keyboard with only three keys: a, b, and c. Two players

bet on whether the pattern "abab" or "aabb" will appear first in the resulting random string. We

analyze who will win this game by computing the expected waiting time until we see each pattern.

We solve the prolem by introducing a casino. At each round, a random letter is drawn. A person

can bet m on the random outcome. He/she will receive 3 m if the guess is correct, but will lose all

of the bet otherwise. This makes it a fair game.

We first consider the pattern "abab"We introduce a team of gamblers, who enter the game

one at a time. Gambler i bets $1 at time i that the i -th letter is the character "a"This gambler will

receive $3 if the letter "a" is drawn at time i but will leave otherwise. Then he will bet $3 to bet

the next drawn letter is "b"He will receive $9 if the letter at time i+ 1 is indeed "b" but will leave

otherwise. The gambler will then bet $27 at time i+ 2 for the letter "a" and $81 at time i+ 3 for

t h el e t t e r" b " H ew i l ll e a v e t h eg a m ew h e n e v e rh el o s e s .

Assume that there are infinitely many potential gamblers waiting outside, and at each time, a

new gambler enters the game. We stop the game when a gambler has just won four times in a row.

This forms a stopped martingale satisfying the third condition in Theorem 13.18. Suppose the game

is stopped at time T At this time, exactly two gamblers are still in the game. The gambler who

entered at timeT- 1 wins $9. - $1 = $8, and the one who entered at time T- 3 wins $81. - $1=$80.

All other gamblers lose $1. By Theorem 13.18, the expected gain of the whole gambling team is 0,

0= E [X0]= E[XT ]= E[(T- 2 )(- $1)+ ( $81- $1 )+ ( $9- $1 )]= $90- $ E[T] .

Therefore, the expected waiting time for the string "abab" is 90.

If we change the pattern to "aabb", only one gambler is in the game when we first see the

pattern. A similar calculation shows that the expected waiting time is 81.

0= E [X0]= E[XT ]= E[(T- 1 )(- $1)+ $81 - $1 ]= $81- $ E[T] .

We simulate this experiment by the following Python program.

from random import choice

A = ['a','b','c'] # alphabet = {a,b,c}

def monkey(alphabet,pattern):

random_string = "" # initialize to empty string

while True:

random_string += choice(alphabet) # randomly draw a character

if random_string[-len(pattern):] == pattern:

break # break if we see the pattern at the end

return len(random_string) # return the length of the random string

n = 10000 # run the experiment n times

waiting_time1 = sum([monkey(A,'abab') for _ in range(n)])/n

waiting_time2 = sum([monkey(A,'aabb') for _ in range(n)])/n

print(f"Average waiting time for pattern abab = {waiting_time1}")

print(f"Average waiting time for pattern aabb = {waiting_time2}")

The function monkey generates a random string that ends with the given pattern and returns

the length of the random string. We run 10,000 experiments for each of the two patterns "abab"

and "aabb" and compute the average length of the random string. Note that the only difference

between the two experiments is the ending patterns. A sample output of the program is

Average waiting time for pattern abab = 89.1899

Average waiting time for pattern aabb = 81.2196

The simulation results are in accordance with the theoretic analysis.
